\documentclass[11pt]{article}

\usepackage[T1]{fontenc}
\usepackage{lmodern}
\usepackage{amsmath,amssymb,amsthm}
\usepackage{booktabs}
\usepackage[margin=1in]{geometry}
\usepackage[hidelinks]{hyperref}
\usepackage{microtype}

\hypersetup{
  pdftitle={No Reversible Gray Representative in a Specified Orbit of a Length-34 Quaternary Legendre Pair},
  pdfauthor={Xiaojian Zeng}
}

\newtheorem{theorem}{Theorem}
\newtheorem{lemma}[theorem]{Lemma}
\newtheorem{proposition}[theorem]{Proposition}

\newcommand{\Z}{\mathbb Z}
\newcommand{\PAF}{\operatorname{PAF}}

\title{No Reversible Gray Representative in a Specified Orbit\\
of a Length-$34$ Quaternary Legendre Pair}
\author{Xiaojian Zeng\\
PBC School of Finance, Tsinghua University, Beijing 100190, China\\
\texttt{xiaojianzeng@foxmail.com}}
\date{}

\begin{document}

\maketitle

\begin{abstract}
Lebedev exhibited an explicit quaternary Legendre pair of length $34$ and
observed that its printed orientation is not reversible under the Gray map,
while leaving an orbit argument unperformed.  We decide a precisely specified
finite-orbit question for that pair.  The allowed operations are independent
fourth-root phases and cyclic shifts, a common unit decimation, simultaneous
complex conjugation, and exchange of the two words.  No member of the resulting
set has four reversible Gray rows.  The proof reduces every permitted transform
of one fixed word to reflection about one of $17$ even cyclic axes, and a
literal mismatch excludes each axis.  A separate exact census of all
$1{,}183{,}744$ parameter tuples corroborates the paper proof.  The conclusion
is only about the displayed set of operations: it neither identifies a full
quaternary-Legendre-pair equivalence group nor excludes reversible
representatives under operations outside that set.
\end{abstract}

\section{Context and the specified orbit}

For a complex word $X=(X_j)_{j\in\Z_n}$, write
\[
  \PAF(X,h)=\sum_{j\in\Z_n}X_j\overline{X_{j+h}}.
\]
A quaternary Legendre pair (qLP) of length $n$ is a pair
$A,B\in\{1,i,-1,-i\}^n$ such that
\[
  \PAF(A,h)+\PAF(B,h)=-2\qquad(0\ne h\in\Z_n).
  \tag{1}
\]
Kotsireas and Winterhof introduced qLPs and related them to quaternary and
binary Hadamard matrices~\cite{KotsireasWinterhof2024}.  The explicit pair
studied below already supplies the existence needed here; our question concerns
only its reversible representatives under specified operations.

Lebedev~\cite{Lebedev2026} printed the following length-$34$ pair.  We encode
$1,i,-1,-i$ by the digits $0,1,2,3$, respectively, and fix
\[
\begin{split}
 A&=2202202000022020222202000220200020,\\
 B&=1133303133112311101113211331303331.
\end{split}
\tag{2}
\]
Lebedev checked that the displayed orientation is not reversible and noted
that an orbit argument would be needed for a conclusion about equivalence
classes.  We now freeze the particular set of operations to be tested.

All indices henceforth lie in $\Z_{34}$.  For $r\in\Z_{34}$ and
$k\in\Z_{34}^{\times}$, define
\[
 (c_rX)_j=X_{j-r},\qquad (d_kX)_j=X_{jk^{-1}}.
 \tag{3}
\]
Let $s(X,Y)=(Y,X)$, and put $X^0=X$ and $X^1=\overline X$.  The set studied
in this note is exactly
\[
 \Gamma_{34}(A,B)=
 \left\{
 s^e\!\left(\alpha c_r d_k(A^\varepsilon),
             \beta c_t d_k(B^\varepsilon)\right):
 \begin{array}{l}
 \alpha,\beta\in\{1,i,-1,-i\},\quad r,t\in\Z_{34},\\
 k\in\Z_{34}^{\times},\quad \varepsilon,e\in\{0,1\}
 \end{array}
 \right\}.
 \tag{4}
\]
The phrase ``specified orbit'' below is shorthand for the displayed finite
set (4).  In particular, it is not a claim that these are all qLP-preserving
operations.

For comparison, the binary Legendre-pair equivalence studied by Bulutoglu,
Baczkowski, and Yauney~\cite{BulutogluBaczkowskiYauney2026} includes formulas
with both common and opposite decimations.  Equation~(4) deliberately uses
only a common decimation; the opposite-decimation alternative is not imported
from the binary setting without a qLP argument.

\begin{proposition}
Every pair in $\Gamma_{34}(A,B)$ satisfies the qLP equations (1).
\end{proposition}

\begin{proof}
For a fourth-root phase $\alpha$, direct reindexing gives
\[
\begin{aligned}
 \PAF(\alpha X,h)&=\PAF(X,h),&
 \PAF(c_rX,h)&=\PAF(X,h),\\
 \PAF(d_kX,h)&=\PAF(X,hk^{-1}),&
 \PAF(\overline X,h)&=\overline{\PAF(X,h)}.
\end{aligned}
\tag{5}
\]
Independent phases and shifts preserve the two summands separately.  A
common decimation permutes the nonzero lags, simultaneous conjugation
conjugates their sum (which is the real number $-2$), and exchange swaps the
summands.  Hence (1) is preserved by every operation in (4).
\end{proof}

\section{The reflection obstruction}

For $p,q\in\{\pm1\}^{34}$, the Gray decoding used here is
\[
 X=\frac{p+q}{2}+i\frac{p-q}{2}.
 \tag{6}
\]
A binary row $x$ is \emph{reversible} when
$x_j=x_{-j}$ for every $j\in\Z_{34}$.

\begin{lemma}[Gray reversibility]
The two Gray rows of a quaternary word $X$ are both reversible if and only if
\[
 X_j=X_{-j}\qquad(j\in\Z_{34}).
 \tag{7}
\]
\end{lemma}

\begin{proof}
On $1,i,-1,-i$, the Gray map has respective values
$(1,1),(1,-1),(-1,-1),(-1,1)$.  It is therefore bijective onto
$\{\pm1\}^2$.  Equality of both Gray coordinates at $j$ and $-j$ is
equivalent to equality of the quaternary symbols there.
\end{proof}

\begin{theorem}
No ordered pair in $\Gamma_{34}(A,B)$ has all four Gray rows reversible.
Equivalently, the existential assertion for the specified set (4) is false.
\end{theorem}

\begin{proof}
Every member of (4), whether exchanged or not, contains in one coordinate a
word
\[
 Z=\alpha c_r d_k(A^\varepsilon).
 \tag{8}
\]
Suppose its two Gray rows were reversible.  By the lemma, $Z_j=Z_{-j}$ for
all $j$.  Cancelling the nonzero phase and, if $\varepsilon=1$, conjugating
both sides yields
\[
 A_{(j-r)k^{-1}}=A_{(-j-r)k^{-1}}\qquad(j\in\Z_{34}).
 \tag{9}
\]
Set $u=(j-r)k^{-1}$.  As $j$ varies, so does $u$, and (9) becomes
\[
 A_u=A_{a-u}\qquad(u\in\Z_{34}),
 \qquad a=-2rk^{-1}\pmod {34}.
 \tag{10}
\]
The residue $a$ is even.  Table~\ref{tab:axes} lists, for every possible even
$a$, an index $u$ at which (10) fails.  Its entries are read directly from
the literal word $A$ in (2).

\begin{table}[ht]
\centering
\caption{A mismatch for every possible reflection axis.}
\label{tab:axes}
\begin{tabular}{rrrrr}
\toprule
$a$ & $u$ & $A_u$ & $a-u\pmod {34}$ & $A_{a-u}$\\
\midrule
 0& 1&2&33&0\\
 2& 0&2& 2&0\\
 4&10&0&28&2\\
 6& 1&2& 5&0\\
 8& 0&2& 8&0\\
10& 0&2&10&0\\
12& 3&2& 9&0\\
14& 1&2&13&0\\
16& 1&2&15&0\\
18& 2&0&16&2\\
20& 0&2&20&0\\
22& 0&2&22&0\\
24& 0&2&24&0\\
26& 3&2&23&0\\
28& 1&2&27&0\\
30& 0&2&30&0\\
32& 1&2&31&0\\
\bottomrule
\end{tabular}
\end{table}

Thus no permitted transform of $A$ has two reversible Gray rows.  The choices
involving $B$ cannot repair that coordinate, and exchange merely changes its
position.  Consequently no pair in (4) has four reversible Gray rows.
\end{proof}

\section{Exact census and scope of the result}

The preceding proof needs only the $17$ literal comparisons in
Table~\ref{tab:axes}.  An exact exhaustive computation was retained as an
independent corroboration.  The parameter domain in (4) has
\[
 4^2\cdot34^2\cdot16\cdot2\cdot2=1{,}183{,}744
 \tag{11}
\]
tuples.  The primary census represents fourth roots by exponents modulo $4$,
enumerates all tuples, and tests the four reflection identities.  It obtains
$1{,}183{,}744$ distinct ordered images and no fully reversible image.  For
each image it stores the first failing Gray row and the least failing index.

A second verifier imports none of the census implementation.  It reconstructs
the parameter domain by separately composed literal-string operations,
rebuilds the ordered images, and recomputes every stored failure witness.  It
reproduces all counts.  This separation between a short mathematical proof and
reconstructible finite evidence is also used in other bounded exact
classifications; for example, Zeng et al.~\cite{ZengEtAl2026Finite} use
symmetry-reduced exhaustive enumeration and explicitly distinguish their
finite-library conclusions from the corresponding unrestricted questions.

The present conclusion has the same kind of strict boundary.  It settles
reversibility only inside (4), for the single pair (2).  It says nothing about
other length-$34$ qLPs and nothing about a larger equivalence relation.  In
particular, the result should not be read as showing that the pair has no
reversible representative under every conceivable qLP-preserving operation.

\paragraph{Data and computational disclosure.}
The exact census, the independent verifier, their machine-readable summaries,
and the complete failure-witness archive accompany the paper.  Both programs
use exact integer and string operations; no random choices, floating-point
tolerances, or external solver enter their acceptance conditions.  From the
project root they are run as
\begin{verbatim}
python3 evidence/orbit_census.py
python3 evidence/verify_orbit_census.py
\end{verbatim}
The exhaustive computation is corroborative: Theorem~3 follows from the
displayed proof without executing either program.

OpenAI Codex assisted with literature search, exact-code preparation and
checking, adversarial organization of the computational audit, and manuscript
preparation.  Language-model output was not used as mathematical evidence.
The theorem-bearing argument is given above, and the retained computations are
reconstructible by the independent verifier.  No proof assistant was used.

\raggedright
\bibliographystyle{abbrvurl}
\bibliography{references}

\end{document}
